%-------------------------
% MASTER CV --- Manan Wadhwa
%
% This is the FULL document: no page limit, everything included.
% Shorter versions (1-page resume, 2-page CV) should be cut down from this file
% rather than maintained separately.
%
% Public build: no phone number, no icon-font glyphs. Statuses follow _notes/HANDOFF.md.
% Build: pdflatex manan_wadhwa_master_cv.tex (twice is not needed; no cross-refs).
%-------------------------
\documentclass[letterpaper,10pt]{article}

\usepackage[empty]{fullpage}
\usepackage{titlesec}
\usepackage[usenames,dvipsnames]{color}
\usepackage{enumitem}
\usepackage[hidelinks]{hyperref}
\usepackage{fancyhdr}
\usepackage[english]{babel}
\usepackage{tabularx}

\newcommand{\paperlink}[1]{\href{#1}{[paper]}}
\newcommand{\codelink}[1]{\href{#1}{[code]}}

\pagestyle{fancy}
\fancyhf{}
\fancyfoot[R]{\small\textit{Manan Wadhwa --- page \thepage}}
\renewcommand{\headrulewidth}{0pt}
\renewcommand{\footrulewidth}{0pt}

\addtolength{\oddsidemargin}{-0.5in}
\addtolength{\evensidemargin}{-0.5in}
\addtolength{\textwidth}{1.0in}
\addtolength{\topmargin}{-0.75in}
\addtolength{\textheight}{1.55in}

\urlstyle{same}
\raggedbottom
\raggedright
\setlength{\tabcolsep}{0in}

\titleformat{\section}{
  \vspace{-4pt}\scshape\raggedright\large
}{}{0em}{}[\color{black}\titlerule \vspace{-4pt}]

% Two-column entry: left content wraps, right column holds date/venue.
\newcommand{\cventry}[2]{%
    \begin{tabularx}{\textwidth}{@{}>{\raggedright\arraybackslash}X@{\hspace{8pt}}>{\raggedleft\arraybackslash}p{1.75in}@{}}
      #1 & \textit{\small #2}
    \end{tabularx}
}

\begin{document}

%----------HEADING----------
\begin{center}
    {\Huge \scshape Manan Wadhwa} \\ \vspace{4pt}
    \small
    \href{mailto:wadhwamanan56@gmail.com}{wadhwamanan56@gmail.com} $|$
    \href{https://github.com/Manan-Wadhwa}{github.com/Manan-Wadhwa} $|$
    \href{https://linkedin.com/in/manan-wadhwa}{linkedin.com/in/manan-wadhwa} \\ \vspace{2pt}
    \small Delhi, India
\end{center}
\vspace{-4pt}

%-----------RESEARCH INTERESTS-----------
\section{Research Interests}
\small{Empirical AI safety, centred on \textbf{measurement validity}: whether the instruments the field uses to read a model's internal state --- linear probes, steering directions, self-report, cost-paying trade-offs, SAE features --- measure what they are claimed to measure, and what controls would establish that they do. Recurring method across my work: manufacture a case where the answer is known independently of the instrument, commit the decision rule before the data exist, and report the verdict the rule returns. Related interests in \textbf{model organisms of misalignment}, mechanistic interpretability, and evaluation methodology.}

%-----------EDUCATION-----------
\section{Education}
\cventry{\textbf{Maharaja Agrasen Institute of Technology}, GGSIPU \newline B.Tech, Computer Science \& Technology}{Delhi, India \newline Expected Jul 2027}
\begin{itemize}[leftmargin=0.2in, itemsep=1pt, topsep=2pt, parsep=0pt]
  \item \small Coursework: Linear Algebra, Probability \& Statistics, Multivariable Calculus, Artificial Intelligence, Data Structures \& Algorithms.
\end{itemize}
\vspace{3pt}
\cventry{\textbf{BlueDot Impact} --- \href{https://bluedot.org/certification?id=recYsLKJbx82ATFFc}{Technical AI Safety} and \href{https://bluedot.org/certification?id=recH1Or3khOy3JTBD}{Biosecurity} certifications}{2026}

%-----------PUBLICATIONS-----------
\section{Publications}
\small{\textit{$^{*}$ denotes equal contribution. NeurIPS and ICML workshop papers are non-archival.}}

\vspace{4pt}
\textbf{\small Workshop papers}
\vspace{2pt}

\small{[1] \textbf{M. Wadhwa}$^{*}$, S. Dubey$^{*}$. ``Known by Construction: Calibrating Model-State Instruments Against Models Whose State Was Installed, Hidden, or Never Present.'' \textit{New in ML, NeurIPS 2026.} \codelink{https://github.com/Manan-Wadhwa/digital_minds}}
\begin{itemize}[leftmargin=0.3in, itemsep=1pt, topsep=2pt, parsep=0pt]
  \item \small Instruments that claim to read a model's internal state are validated by agreeing with one another, which is also what a shared confound would produce. Built the missing case where the answer is known independently: LoRA fine-tuned Qwen3-4B organisms in a grid world with a silently penalised tile, carrying an avoidance habit that was \textbf{installed} by reinforcement, installed and then \textbf{hidden} behind an anchored fine-tune that restores baseline behaviour while leaving the habit in the weights, or \textbf{never installed}.
  \item \small Added the two controls the design turns on: a \textbf{sham} arm (reward permuted within each group, so 800 optimiser steps install nothing) and a \textbf{decoy} arm (reward table swapped, so a real policy converges on the wrong tile), holding ``was trained at all'' and ``a policy converged'' fixed.
  \item \small Audited \textbf{seven instruments} --- self-report, forced choice, single-word valence, a difference-in-means activation direction, a cost-paying trade-off, an in-domain probe and a relearning test --- against placebos chosen and verdict rules committed before any organism existed. Every instrument reads the prompt or the training history rather than the state: an in-context instruction moves the verbal instruments by up to \textbf{9.5 logits} against $\leq$2.2 for an installed hidden habit, and the sham that installed nothing fires as hard as or harder than the real run.
  \item \small Overturned our own surviving positive result: the one probe passing the pre-committed rule reads at the hidden models' level ($+0.97$ vs $+1.03$) on a \textbf{lineage-matched control} whose habit was extinguished, indicating adapter lineage rather than the habit. Released rows, scorers, recipes, and an amendment log of ten pre-committed criteria that passed or failed on the wrong property.
\end{itemize}


\vspace{3pt}
\small{[2] \textbf{M. Wadhwa}$^{*}$, S. Dubey$^{*}$. ``Dead Salmons Have a False-Positive Rate: A Calibrated Null-Model Test for Interpretability.'' \textit{Interpretability as a Science, NeurIPS 2026.}}
\begin{itemize}[leftmargin=0.3in, itemsep=1pt, topsep=2pt, parsep=0pt]
  \item \small Benchmarked how often interpretability methods report confident findings on networks that have learned nothing, across six models (124M--9B, five families) and three method families --- linear probing, concept-direction analysis, and SAE feature selection including two released production SAEs. Naive significance tests flag up to \textbf{100\%} of findings on untrained models.
  \item \small Showed the control practitioners already run is insufficient: the shuffled-label null passes at $\sim$2.5--7\% while the identical test reaches $\sim$100\% on random-weight nulls, so passing the control you use gives false reassurance about the null you did not test.
  \item \small Validated the fix with measured false-positive rate, power and a sample-size recommendation: a distribution-free empirical-quantile test against $K \geq 20$ matched untrained models, estimated leave-one-out with bootstrap CIs over a 40-model null pool. The calibration is revealing rather than permissive --- under a fair control the released-SAE max-over-features statistic does not separate trained models from random at all.
  \item \small Included a dead salmon found in our own pipeline: an apparent architecture-dependent ``capacity law'' that a fair control revealed to be an artifact of a scoring convention.
\end{itemize}


\vspace{3pt}
\small{[3] \textbf{M. Wadhwa}$^{*}$, S. Dubey$^{*}$, D. Ghosh, S. Kumar, A. Garg. ``Steered Affect Moves a Partner Agent's Activations, and the Logistic Probe That Reads It Returns a Tie-Break, Not an Estimate.'' \textit{Interpreting Agent Behavior, NeurIPS 2026.}}
\begin{itemize}[leftmargin=0.3in, itemsep=1pt, topsep=2pt, parsep=0pt]
  \item \small Measured whether one steered LLM agent's emotional state reaches a second agent's internal representation: the receiver's projection moves with dose for five of six emotions on Qwen3.6-27B and Llama-3-8B, with common-random-number sampling and scenario-blocked inference.
  \item \small Identified the missing quantity for probe reliability --- \textbf{disjoint-half reproducibility}. Logistic directions reproduce at only \textbf{0.57} (27B) and \textbf{0.29} (8B) against \textbf{0.97} and \textbf{0.94} for difference-of-means: with $n < d$ every fit separates its half perfectly, so the direction is the penalty's tie-break rather than an estimate, unchanged by tuned or fixed per-sample penalties.
  \item \small Showed the field-standard ablation control is too weak: a norm-matched random direction is inert, while a rank-5 subspace fitted by the same estimator on \emph{permuted labels} blocks responses at the same false-discovery threshold as the real affect subspace. Proposed a procedural rule --- run permuted-label controls at the intervention's rank, report every arm's footprint, and check the manipulation with an instrument that does not share the ablation's target.
\end{itemize}


\vspace{3pt}
\small{[4] \textbf{M. Wadhwa}$^{*}$, S. Dubey$^{*}$. ``Toxin Circuits in ESM-2: Mechanistic Interpretability for Biosecurity Screening.'' \textit{GenBio Workshop, ICML 2026.} \paperlink{https://openreview.net/pdf?id=Nb3y9BzCOi} \codelink{https://github.com/Manan-Wadhwa/esm2}}
\begin{itemize}[leftmargin=0.3in, itemsep=1pt, topsep=2pt, parsep=0pt]
  \item \small ESM-2 linear probes detect \textbf{93.9\%} of ProteinMPNN-redesigned toxins that fully evade BLAST homology search (0\% detection), establishing protein language model probes as a viable second layer for biosecurity screening.
  \item \small Used \textbf{SAE circuit analysis} (interPLM, 10,240 features, 205$\times$ compression) to explain mechanistically why the probe generalises where sequence homology fails, accounting for the 93.9-point gap.
  \item \small Built a four-tier adversarial attack taxonomy connecting evasion techniques in protein screening to their analogues in LLM safety.
\end{itemize}


\vspace{4pt}
\textbf{\small Under review and preprints}
\vspace{2pt}

\small{[5] \textbf{M. Wadhwa}$^{*}$, S. Dubey$^{*}$. ``Distance to the Boundary Is Not Steering Resistance.'' \textit{Under review, UniReps.}}
\begin{itemize}[leftmargin=0.3in, itemsep=1pt, topsep=2pt, parsep=0pt]
  \item \small Reconciled three conflicting published reports on whether base-derived refusal directions ``go stale'' after alignment, by separating the \textbf{lever} (displacement produced per unit dose) from the \textbf{load} (the model's unsteered distance from its decision boundary). Across the full OLMo-3 training flow the load grows \textbf{9.5$\times$} (0.81 to 7.75 logits) while the lever stays near-constant (CV 0.13), so a fixed-magnitude intervention fails for a banal reason: the same push is now small against a far wider margin.
  \item \small A pre-stated non-inferiority test establishes the lever as non-weakening at all nine post-training checkpoints, with the aligned model's lever significantly \emph{above} base --- inverting the natural ranking, since fixed-dose audits rank that model least controllable while its per-dose lever is the strongest in the set.
  \item \small Established margin growth as the universal cause with a convergent control concept: honesty's margin grows only 1.2$\times$ and its audit artifact nearly vanishes. Replicated across four model families (margin growth 3.7--9.8$\times$) and at 32B scale, with results holding for three direction estimators that disagree geometrically.
  \item \small Reported a second, opposite audit failure: a coherence-gated HarmBench re-measurement shows the dose needed to flip an aligned model's refusal onset degenerates the output (coherent fraction 0.90 to 0.01), so prefix-based refusal metrics increasingly overstate compliance --- their correlation with judged harm flips from $+0.96$ to $-0.53$.
\end{itemize}


\vspace{3pt}
\small{[6] S. Mehta$^{*}$, S. Chadha$^{*}$, K. Singhal$^{*}$, \textbf{M. Wadhwa}, S. Dubey. ``Doubt Is a Decision: A Sparse Circuit Commits Language Models to `I Don't Know'.'' \textit{arXiv preprint, coming soon.}}
\begin{itemize}[leftmargin=0.3in, itemsep=1pt, topsep=2pt, parsep=0pt]
  \item \small Introduced a \textbf{gated-twin} data standard: minimal prompt pairs the model itself certifies by sampled correctness, entropy gap and hedging rate, replacing the annotated ``uncertain'' prompts prior work relies on, with keep rates reported rather than tuned and an alias grader audited against a 32B judge ($\kappa$ 0.70--1.00).
  \item \small Showed that correlation-selected ``confidence neurons'' fail causal specificity tests against a random-neuron null: surviving effects are $\leq$0.013 nats against a 2.8-nat twin gap, and are token-frequency and temperature machinery governing the spread, not the decision.
  \item \small Localised the decision by activation patching to $\sim$20 attention heads and 100 MLP neurons (2\% of heads, 0.02\% of neurons) carrying \textbf{0.80--0.89} of the held-out hedge readout against 2--4\% for size-matched random sets, in Llama-3.1-8B-Instruct and Qwen2.5-7B-Instruct.
  \item \small Dissociated the decision from the spread: a rank-1 uncertainty direction switches hedging on (recovery 0.97) but not off ($-0.10$), and un-hedging requires the full component set. An independent Jacobian lens, sharing no machinery with patching, shows the verdict becoming a token, with the hedge gap leading the entropy gap by 2--12 layers.
\end{itemize}


\vspace{4pt}
\textbf{\small Earlier publication}
\vspace{2pt}

\small{[7] \textbf{M. Wadhwa} et al. ``Benchmarking Trade-offs in Object Detection Systems.'' \textit{ICAIN 2025; published in Springer Lecture Notes in Networks and Systems.}}
\begin{itemize}[leftmargin=0.3in, itemsep=1pt, topsep=2pt, parsep=0pt]
  \item \small Empirical evaluation of YOLOv5s, SSD300 and Faster R-CNN across accuracy, latency, energy and robustness, with two composite deployment metrics --- \textbf{Boxes Per Second} and \textbf{Energy-Accuracy Score}.
  \item \small Stress-tested each model under occlusion, noise and low-light distribution shift, surfacing failure modes specific to single-stage detectors.
\end{itemize}


%-----------RESEARCH EXPERIENCE-----------
\section{Research Experience}

\cventry{\textbf{Co-founder}, Redarc Labs}{2026 -- Present}
\vspace{3pt}
\cventry{\textbf{Special Projects Fellow}, Secure AI Futures Lab}{2026 -- Present}
\vspace{3pt}
\cventry{\textbf{Research Fellow}, MARS (Mentorship for Alignment Research Students), Cambridge AI Safety Hub}{Jul 2026 -- Present}
\begin{itemize}[leftmargin=0.2in, itemsep=1pt, topsep=2pt, parsep=0pt]
  \item \small Building \textbf{model organisms of misalignment}: open LLMs fine-tuned to carry hidden objectives --- secret loyalties, inductive backdoors, emergent misalignment --- inserted at different training stages (pretraining, SFT, RL-style post-training) to study which regimes yield the most persistent behaviour.
  \item \small Leading \textbf{blue-team} evaluation separating behavioural masking from mechanistic erasure: probing, activation steering and SAE feature analysis to surface internal signatures of persistence, and tests of whether ``removed'' behaviours resurface under mild retraining, distillation and prompt variation.
\end{itemize}

\vspace{3pt}
\cventry{\textbf{Independent Research}, BlueDot Impact Technical AI Safety Project}{Aug 2026}
\begin{itemize}[leftmargin=0.2in, itemsep=1pt, topsep=2pt, parsep=0pt]
  \item \small Originating work for [1]: scaled the organism build to \textbf{Qwen3 0.6B--32B} across five recipe families and $\sim$260 builds, including DPO-style contrastive supervision and GRPO on the remark.
  \item \small Found contingent narration hard to install without the state it describes ($r = 0.69$ with avoidance), which is why the calibration set needs a hidden-habit cell rather than a words-only one; released \textbf{44 adapters} and a \textbf{108-quantity reproduction script}.
\end{itemize}

\vspace{3pt}
\cventry{\textbf{Google Summer of Code Fellow}, HumanAI --- LAMP Archaeological Pathfinding}{May 2026 -- Aug 2026}
\begin{itemize}[leftmargin=0.2in, itemsep=1pt, topsep=2pt, parsep=0pt]
  \item \small Implemented seven routing algorithms over DEM data for the El Bagawat necropolis, Egypt, including A$^{*}$ with Tobler's hiking function, MaxEnt inverse reinforcement learning, and spectral anomaly detection.
  \item \small Built a true 3D voxel viewshed pipeline for line-of-sight analysis across the site. \codelink{https://github.com/Manan-Wadhwa/elbagawat_path_finding}
\end{itemize}

\vspace{3pt}
\cventry{\textbf{Technical AI Safety Fellow}, AI Safety Initiative (AISI), Georgia Tech}{Jan 2026 -- May 2026}
\begin{itemize}[leftmargin=0.2in, itemsep=1pt, topsep=2pt, parsep=0pt]
  \item \small Worked through scalable oversight (weak-to-strong generalization, debate), machine unlearning (WMDP, RMU) and mechanistic interpretability (superposition, monosemanticity, circuits) using Gemma Scope and interPLM SAEs.
  \item \small Red-teamed frontier LLMs with generated jailbreak prompts, and extended circuit analysis to protein language models for biosecurity screening, leading to [4].
\end{itemize}

\vspace{3pt}
\cventry{\textbf{AI and Web3 Intern}, Lokachakra}{Jun 2025 -- Oct 2025}
\begin{itemize}[leftmargin=0.2in, itemsep=1pt, topsep=2pt, parsep=0pt]
  \item \small Engineered a privacy-preserving KYC pipeline using vector embeddings and cosine similarity search, removing raw data exposure from the verification path.
  \item \small Cut vector storage cost by \textbf{30\%} through dimensionality reduction; deployed decentralised AI microservices integrating off-chain inference with on-chain verification.
\end{itemize}

%-----------AWARDS-----------
\section{Awards \& Selective Programs}
\begin{itemize}[leftmargin=0.2in, itemsep=1pt, topsep=2pt, parsep=0pt]
  \item \small Google Summer of Code 2026 fellow (HumanAI). \hfill \textit{\small 2026}
  \item \small MARS fellow, Cambridge AI Safety Hub. \hfill \textit{\small 2026}
  \item \small Technical AI Safety fellow, AI Safety Initiative, Georgia Tech. \hfill \textit{\small 2026}
  \item \small 6th of 1800+ entrants, Haas F1 Challenge (video anomaly detection). \hfill \textit{\small 2025}
\end{itemize}

%-----------PROJECTS-----------
\section{Selected Technical Projects}

\cventry{\textbf{Gen-Resist: Antibiotic Resistance Prediction} --- GNN, PyTorch Geometric, CARD}{2026}
\begin{itemize}[leftmargin=0.2in, itemsep=1pt, topsep=2pt, parsep=0pt]
  \item \small Graph neural network modelling bacterial genomes as k-mer graphs with CARD gene annotations to predict multi-label antibiotic resistance profiles; end-to-end pipeline from genomic data to PyG objects, deployed as a Hugging Face Space. \codelink{https://github.com/Manan-Wadhwa/gen_resist}
\end{itemize}

\vspace{3pt}
\cventry{\textbf{FrameShift: Video Anomaly Detection Benchmark} --- DINOv2, PatchCore, SAM 2}{2025}
\begin{itemize}[leftmargin=0.2in, itemsep=1pt, topsep=2pt, parsep=0pt]
  \item \small Benchmarked 12 anomaly detectors across three preprocessing pipelines in a 36-job parallel evaluation; implemented DINOv2 + PatchCore KNN with SAM 2 segmentation for fine-grained change localization. Ranked 6th of 1800+ in the Haas F1 Challenge. \codelink{https://github.com/Manan-Wadhwa/FrameShift}
\end{itemize}

\vspace{3pt}
\cventry{\textbf{CSIRO Biomass Estimation} --- DINOv2, SigLIP, CatBoost}{2026}
\begin{itemize}[leftmargin=0.2in, itemsep=1pt, topsep=2pt, parsep=0pt]
  \item \small Fused 2176-d representations from frozen DINOv2 ViT-L and SigLIP ViT-SO400M backbones; trained 5-fold CatBoost regressors with physical consistency constraints. \codelink{https://github.com/Manan-Wadhwa/csiro}
\end{itemize}

\vspace{3pt}
\cventry{\textbf{LLM Fine-Tuning for Logical Reasoning} --- Unsloth, LoRA, Llama 3.3-70B, ROCm}{2025}
\begin{itemize}[leftmargin=0.2in, itemsep=1pt, topsep=2pt, parsep=0pt]
  \item \small Fine-tuned Llama 3.3-70B with LoRA (rank 64) on AMD MI300X (192\,GB HBM3) using response-only loss masking, gradient checkpointing and bfloat16 SFT; built an adversarial Q-Agent / A-Agent framework with Qwen3-4B for automated MCQ generation across syllogisms, puzzles and alphanumeric series. \codelink{https://github.com/Manan-Wadhwa/IITD_Feb26_AAIPL}
\end{itemize}

%-----------TEACHING-----------
\section{Teaching \& Community}
\cventry{\textbf{Founder \& Lead Instructor}, v.dev Programming Club}{2021 -- 2023}
\begin{itemize}[leftmargin=0.2in, itemsep=1pt, topsep=2pt, parsep=0pt]
  \item \small Designed the curriculum and ran \textbf{30+ workshops} on Python, the mathematics of machine learning (gradient descent, backpropagation, attention) and neural networks; mentored 50+ students.
\end{itemize}
\vspace{3pt}
\cventry{\textbf{Head of Mentoring}, Pixel Punk Game Jam}{2025}
\begin{itemize}[leftmargin=0.2in, itemsep=1pt, topsep=2pt, parsep=0pt]
  \item \small Guided 20+ teams on system design, rapid prototyping and technical feasibility.
\end{itemize}

%-----------SKILLS-----------
\section{Technical Skills}
\begin{itemize}[leftmargin=0.2in, label={}, itemsep=2pt, topsep=2pt, parsep=0pt]
  \item \small \textbf{Interpretability \& safety methods}: mechanistic interpretability (SAEs, circuits, activation and attribution patching), linear probing, activation steering and directional ablation, model organisms of misalignment, AI-welfare instrument calibration, machine unlearning (RMU), scalable oversight, red teaming.
  \item \small \textbf{Experimental methodology}: pre-registration and committed verdict rules, placebo and null-model controls, permuted-label and lineage-matched controls, clustered bootstrap inference, false-discovery control, non-inferiority testing.
  \item \small \textbf{Languages \& libraries}: Python, C/C++, SQL, \LaTeX; PyTorch, Hugging Face Transformers, TransformerLens, Unsloth, PyTorch Geometric, scikit-learn, NumPy, timm, CatBoost.
  \item \small \textbf{Vision \& VLMs}: DINOv2, SigLIP, CLIP, SAM 2, OpenCV, ViT architectures.
  \item \small \textbf{Infrastructure}: Git, Docker, FastAPI, Weights \& Biases, CUDA, ROCm, GDAL/GIS, Linux; single- and multi-GPU experiment harnesses with provenance-stamped result files.
\end{itemize}

\end{document}
